% Minimal preamble for Open Logic Project sources (pinned commit 1e960beff9ed7835bf3e3f1335e21af3439cd107).
%
% OLP's own master preamble cannot run on the project host: TeX Live 2025 there lacks
% prooftree.sty, ntheorem.sty, extarrows.sty, mathdots.sty, isomath.sty and subfiles.sty.
% Following the project's own rule ("if the author preamble needs packages the toolchain lacks,
% do not invent macros -- use a minimal working preamble"), this file declares, verbatim in
% spirit and numbering behaviour, exactly the theorem-like environments OLP defines in
% sty/open-logic-book-envs.sty (thm/lem/prop/cor/defn/ex/prob/rem/note/case/axiom/conv),
% plus the shorthand accents OLP uses (`!!a`, `!!b`, `!!{...}`).
%
% Author text is untouched; only the packaging is ours.
\documentclass[10pt,a4paper]{article}

\usepackage{amsmath}
\usepackage{amssymb}
\usepackage{amsthm}
\usepackage{graphicx}
\usepackage{hyperref}
\usepackage{enumitem}
\usepackage[framemethod=default]{mdframed}

% ---- OLP shorthand accents (from sty/open-logic.sty) ----
% NB: \! is taken by LaTeX (negative thin space), so we only define the !!a / !!b families
% and the \olup helper used by them.
\newcommand*{\olup}[1]{\MakeUppercase{#1}}
\def\!!!a#1{\olup{#1}}
\def\!!a#1{\olup{#1}}
\def\!!b#1{\olup{#1}}

% ---- theorem-like environments, same family numbering as OLP ----
\theoremstyle{plain}
\newtheorem{thm}{Theorem}[section]
\newtheorem{lem}[thm]{Lemma}
\newtheorem{prop}[thm]{Proposition}
\newtheorem{cor}[thm]{Corollary}
\newtheorem{axiom}{Axiom}

\theoremstyle{definition}
\newtheorem{defn}[thm]{Definition}
\newtheorem{ex}[thm]{Example}
\newtheorem{prob}{Problem}[section]
\newtheorem{probd}{Problem}
\newtheorem{conv}[thm]{Convention}

\theoremstyle{remark}
\newtheorem{rem}[thm]{Remark}
\newtheorem{note}[thm]{Note}
\newtheorem{case}[thm]{Case}

% ---- environments OLP defines as wrappers ----
\newenvironment{probtag}[1]{\begin{prob}}{\end{prob}}
\newenvironment{explain}{}{}
\newenvironment{explain*}{}{}
\newenvironment{editorial}{}{}
\newenvironment{digress}{}{}
\newenvironment{digress*}{}{}
\newenvironment{history}{}{}
\newenvironment{reference}{}{}
\newcommand*{\olfileid}[3]{}
\newcommand*{\olsection}[1]{\section*{#1}}
\newcommand*{\olchapter}[3]{}
\newcommand*{\olpart}[2]{}
\newcommand*{\olimport}[1]{}
\newcommand*{\olinclude}[1]{}
\newcommand*{\olpath}[1]{}
\newcommand*{\ollangid}[1]{}
\newcommand*{\problemtag}[1]{}
\newcommand*{\problemauthor}[1]{}
\newcommand*{\ollabel}[1]{\label{#1}}
% \olref takes 0-3 OPTIONAL arguments plus one mandatory label:
%     \olref{lab}  |  \olref[from]{lab}  |  \olref[from][to]{lab}  |  \olref[a][b][c]{lab}
% A single-argument definition silently swallows "[sth]" and emits a broken reference
% (measured on the host: the text came out as "??iso]x"). The chain below consumes the
% brackets and cross-references only the mandatory argument.
\makeatletter
\newcommand{\olref}{\@ifnextchar[{\olref@i}{\olref@final}}
\def\olref@i[#1]{\@ifnextchar[{\olref@ii}{\olref@final}}
\def\olref@ii[#1]{\@ifnextchar[{\olref@iii}{\olref@final}}
\def\olref@iii[#1]{\olref@final}
\def\olref@final#1{\ref{#1}}
\makeatother
\newcommand*{\olnewpage}{\clearpage}

% ===================== OLP compatibility layer (1) =====================
% Compatibility layer for Open Logic Project sources (pinned commit 1e960beff9ed7835bf3e3f1335e21af3439cd107).
%
% OLP's own style stack (sty/open-logic.sty) needs ntheorem, prooftree, extarrows, mathdots,
% isomath and subfiles; the project's local and host TeX Live installs lack several of those.
% Rather than vendor packages, this file supplies drop-in equivalents for every OLP construct
% that the frozen units actually use. It only rearranges the author's own tokens; no
% mathematical content is added, and nothing here is claimed to be OLP's own file.

\makeatletter

% ---------------- proof trees (OLP's prooftree is a bussproofs derivative) ----------------
% These commands CONSUME their argument and typeset nothing: the bracket labels of OLP's proof
% trees carry rule names. Printing them inline dropped those labels into math mode and produced
% 63 "Missing $" / "allowed only in math mode" errors; swallowing them keeps the brace balance
% and the mode correct. The derivation text between the commands is still printed unchanged.
\providecommand{\AxiomC}[1]{}
\providecommand{\UnaryInfC}[1]{}
\providecommand{\BinaryInfC}[1]{}
\providecommand{\TrinaryInfC}[1]{}
\providecommand{\QuaternaryInfC}[1]{}
\providecommand{\QuinaryInfC}[1]{}
\providecommand{\RightLabel}[1]{}
\providecommand{\LeftLabel}[1]{}
\providecommand{\doubleLine}{}
\providecommand{\noLine}{}
\providecommand{\DisplayProof}{}
\providecommand{\fCenter}{}
\newenvironment{prooftree}{}{}
\newenvironment{ltprooftree}{}{}
\newenvironment{oltableau}{}{}

% ---------------- logical notation ----------------
\providecommand{\lif}{\mathbin{\rightarrow}}
\providecommand{\lforall}{\mathord{\forall}}
\providecommand{\lexists}{\mathord{\exists}}
\providecommand{\ltrue}{\top}
\providecommand{\lfalse}{\bot}
\providecommand{\Ltrue}{\top}
\providecommand{\Lfalse}{\bot}
\providecommand{\Sequent}{\mathrel{\Rightarrow}}
\providecommand{\Proves}{\mathrel{\vdash}}
\providecommand{\Entails}{\mathrel{\models}}
\providecommand{\Subst}[3]{#1[#2/#3]}
\providecommand{\sFmla}[2]{#1(#2)}
\providecommand{\pFmla}[2]{#1(#2)}
\providecommand{\formula}{\relax}
\providecommand{\True}{\top}
\providecommand{\False}{\bot}
\providecommand{\Args}[1]{#1}
\providecommand{\Ts}{\ensuremath{\top}}
\providecommand{\Fi}{\ensuremath{\bot}}
\providecommand{\lfalsum}{\bot}

% ---------------- objects, structures, metatheory ----------------
\providecommand{\Obj}{}
\providecommand{\Struct}{\ensuremath{\mathfrak{A}}}
\providecommand{\Domain}{\ensuremath{|\mathfrak{A}|}}
\providecommand{\Assign}{\ensuremath{s}}
\providecommand{\Sat}[2]{\ensuremath{#1 \models #2}}
\providecommand{\mSat}[3]{\ensuremath{#1 \models #2[#3]}}
\providecommand{\mModel}[2]{\ensuremath{#1 \models #2}}
\providecommand{\Lang}{\ensuremath{\mathcal{L}}}
\providecommand{\Th}{\ensuremath{\mathrm{Th}}}
\providecommand{\Log}{\ensuremath{\mathrm{L}}}
\providecommand{\tuple}[1]{\langle #1 \rangle}
\providecommand{\num}[1]{\underline{#1}}
\providecommand{\cfind}[1]{\varphi_{#1}}
\providecommand{\Nat}{\ensuremath{\mathbb{N}}}
\providecommand{\Mult}{\ensuremath{\mathrm{Mult}}}
\providecommand{\Add}{\ensuremath{\mathrm{Add}}}
\providecommand{\Exp}{\ensuremath{\mathrm{Exp}}}
\providecommand{\Zero}{\ensuremath{0}}
\providecommand{\Succ}{\ensuremath{\mathrm{S}}}
\providecommand{\lambd}{\ensuremath{\lambda}}
\providecommand{\bredpar}{\mathrel{\rightarrow_\beta}}
\providecommand{\bcd}[1]{\ensuremath{\ulcorner #1 \urcorner}}
\providecommand{\Set}[1]{\ensuremath{\{#1\}}}
\providecommand{\indcase}[1]{\par\noindent\emph{Case #1}.\quad}
\providecommand{\indfrm}[1]{\par\noindent\emph{Form #1}.\quad}
\providecommand{\iftag}[1]{#1}
\providecommand{\tagitem}[1]{\item[#1]}
% these five are used with arguments by OLP's proof trees -> swallow the argument
\providecommand{\Elim}[1]{}
\providecommand{\Intro}[1]{}
\providecommand{\LeftR}[1]{}
\providecommand{\RightR}[1]{}
\providecommand{\TAss}[1]{}
\providecommand{\TRule}[1]{}

\makeatother
\providecommand{\olpath}{.}

% ===================== OLP compatibility layer (2) =====================
% --- second compatibility block: notation and shorthands used by the frozen OLP units ---
% Everything here is \providecommand, so anything already defined by LaTeX or the packages wins.
\providecommand{\CutCS}{\mathrm{Cut}}
\providecommand{\LeftSubproofLabel}{\relax}
\providecommand{\RightSubproofLabel}{\relax}
\providecommand{\maeh}{\relax}
\providecommand{\cheight}{\relax}
\providecommand{\ande}{\mathrm{and}_E}
\providecommand{\andi}{\mathrm{and}_I}
\providecommand{\ore}{\mathrm{or}_E}
\providecommand{\ori}{\mathrm{or}_I}
\providecommand{\typeof}{\mathord{:}}
\providecommand{\pValue}[1]{#1}
\providecommand{\Sent}{\mathrm{Sent}}
\providecommand{\cardplus}{\mathbin{+}}
\providecommand{\cardtimes}{\mathbin{\cdot}}
\providecommand{\cardexpo}{\mathbin{\hat{\ }}}
\providecommand{\cardnless}{\not<}
\providecommand{\cardle}{\leq}
\providecommand{\funfromto}[3]{#1\colon #2 \to #3}
\providecommand{\funimage}[1]{\operatorname{im}(#1)}
\providecommand{\fdefined}{\mathord{\downarrow}}
\providecommand{\canonord}{\operatorname{ord}}
\providecommand{\ordtimes}{\mathbin{\cdot}}
\providecommand{\ordeq}[2]{#1 = #2}
\providecommand{\ordneq}[2]{#1 \neq #2}
\providecommand{\ordtype}[1]{\operatorname{ord}(#1)}
\providecommand{\ordsucc}{\mathord{+1}}
\providecommand{\ordunion}{\cup}
\providecommand{\onesphere}{S^1}
\providecommand{\rotationsgroup}{\mathrm{SO}(2)}
\providecommand{\equivrep}[1]{[#1]}
\providecommand{\rlexless}{\prec}
\providecommand{\ZFC}{\mathrm{ZFC}}
\providecommand{\ZF}{\mathrm{ZF}}
\providecommand{\Rat}{\mathbb{Q}}
\providecommand{\Real}{\mathbb{R}}
\providecommand{\PosInt}{\mathbb{Z}^{+}}
\providecommand{\Id}{\mathrm{id}}
\providecommand{\Bin}{\mathrm{Bin}}
\providecommand{\pto}{\rightharpoonup}
\providecommand{\dom}{\operatorname{dom}}
\providecommand{\ran}{\operatorname{ran}}
\providecommand{\len}{\operatorname{len}}
\providecommand{\concat}{\mathbin{{}^\frown}}
\providecommand{\card}[1]{\lvert #1 \rvert}
\providecommand{\Pow}{\mathcal{P}}
\providecommand{\Setabs}[2]{\{\, #1 \mid #2 \,\}}
\providecommand{\comp}{\circ}
\providecommand{\Char}{\chi}
\providecommand{\Prf}{\operatorname{Prf}}
\providecommand{\Prov}{\operatorname{Prov}}
\providecommand{\Proj}{\operatorname{proj}}
\providecommand{\Gn}{\operatorname{gn}}
\providecommand{\gn}{\operatorname{gn}}
\providecommand{\Value}[1]{\operatorname{val}(#1)}
\providecommand{\Deduce}{\mathrm{D}}
\providecommand{\Dual}{\mathrm{d}}
\providecommand{\Ax}{\mathrm{Ax}}
\providecommand{\pSat}[3]{#1 \models #2\,[#3]}
\providecommand{\pAssign}{\sigma}
\providecommand{\bforall}{\forall}
\providecommand{\bexists}{\exists}
\providecommand{\liff}{\mathbin{\leftrightarrow}}
\providecommand{\FV}{\operatorname{FV}}
\providecommand{\UnaryInf}{\ensuremath{\mid}}
\providecommand{\BinaryInf}{\ensuremath{\mid}}
\providecommand{\depth}{\operatorname{depth}}
\providecommand{\oliflabeldef}[3]{#1}
\providecommand{\TMblank}{\sqcup}
\providecommand{\TMstroke}{\mid}
\providecommand{\TMstay}{\mathrm{S}}
\providecommand{\TMright}{\mathrm{R}}
\providecommand{\TMleft}{\mathrm{L}}
\providecommand{\TMendtape}{\triangleright}
\providecommand{\umin}{\mu}
\providecommand{\Sum}{\sum}
% bibliography helpers (natbib): the frozen units cite the shared OLP bibliography
\providecommand{\citep}[1]{\cite{#1}}
\providecommand{\citet}[1]{\cite{#1}}
\providecommand{\citeauthor}[1]{\cite{#1}}
\providecommand{\citeyear}[1]{\cite{#1}}

% ===================== OLP compatibility layer (3) =====================
% --- third compatibility block: remaining notation seen in the frozen OLP units ---
\providecommand{\iso}{\cong}
\providecommand{\nszero}{\mathrm{0}}
\providecommand{\nssucc}{\mathrm{S}}
\providecommand{\nsplus}{\mathbin{+}}
\providecommand{\nssless}{<}
\providecommand{\nsless}{<}
\providecommand{\PL}{\mathrm{PL}}
\providecommand{\Nec}{\mathord{\Box}}
\providecommand{\RK}{\mathrm{RK}}
\providecommand{\ST}{\mathrm{ST}}
\providecommand{\SSubst}[3]{#1[#2/#3]}
\providecommand{\subst}[3]{#1[#2/#3]}
\providecommand{\tagfalse}{\mathrm{F}}
\providecommand{\cutr}{\mathrm{Cut}}
\providecommand{\cardneq}{\neq}
\providecommand{\issucc}[1]{\operatorname{succ}(#1)}
\providecommand{\tup}[1]{\langle #1 \rangle}
\providecommand{\Var}{\mathrm{Var}}
\providecommand{\Bew}{\operatorname{Bew}}
\providecommand{\Con}{\operatorname{Con}}
\providecommand{\Cons}{\operatorname{Cons}}
\providecommand{\pDeduce}{\mathrm{D}}
\providecommand{\psat}{\models}
\providecommand{\den}[1]{\llbracket #1 \rrbracket}
\providecommand{\Sem}[2]{\llbracket #1 \rrbracket_{#2}}
\providecommand{\DomainOf}[1]{\lvert #1 \rvert}
\providecommand{\Form}{\mathrm{Form}}
\providecommand{\Terms}{\mathrm{Term}}
\providecommand{\Atom}{\mathrm{Atom}}
\providecommand{\FVars}[1]{\operatorname{FV}(#1)}
\providecommand{\fv}[1]{\operatorname{fv}(#1)}
\providecommand{\tuplep}[1]{\langle #1 \rangle}
\providecommand{\restrict}[2]{#1 \upharpoonright #2}
\providecommand{\sto}{\rightarrow}
\providecommand{\FmlaSet}[1]{\mathrm{Fmla}_{#1}}
\providecommand{\Thm}{\mathrm{Thm}}
\providecommand{\Eval}{\operatorname{val}}
\providecommand{\blank}{\sqcup}
\providecommand{\blankTape}{\sqcup}
\providecommand{\Consis}{\operatorname{Con}}
\providecommand{\drs}{\mathrm{d}}
\providecommand{\VRule}{\mid}
\providecommand{\Bet}{\mathrm{B}}
\providecommand{\Seg}[1]{\mathrm{Seg}_{#1}}
\providecommand{\Enums}{\mathrm{En}}
\providecommand{\Z}{\mathbb{Z}}
\providecommand{\N}{\mathbb{N}}
\providecommand{\abs}[1]{\lvert #1 \rvert}
\providecommand{\floor}[1]{\lfloor #1 \rfloor}
\providecommand{\ceil}[1]{\lceil #1 \rceil}

% ===================== OLP compatibility layer (4) =====================
% --- fourth compatibility block: harvested from the per-candidate compile logs ---
\providecommand{\beredpar}{\mathrel{\rightarrow_\beta}}
% \bered / \bred are used as infix relations AND with a bracketed annotation, so they need the
% optional-argument chain; \providecommand would leave an argument-less definition in place.
\makeatletter
\def\bered{\@ifnextchar[{\bered@i}{\mathrel{\rightarrow_\beta}}}
\def\bered@i[#1]{\mathrel{\rightarrow_\beta}}
\def\bred{\@ifnextchar[{\bred@i}{\mathrel{\rightarrow_\beta}}}
\def\bred@i[#1]{\mathrel{\rightarrow_\beta}}
\makeatother
\providecommand{\beredone}{\mathrel{=_\beta}}
\providecommand{\bredone}{\mathrel{=_\beta}}
\providecommand{\xredone}{\mathrel{=_\beta}}
\providecommand{\aconvone}{\mathrel{=_\alpha}}
\providecommand{\cardless}{<}
\providecommand{\tagrefs}[1]{#1}
\providecommand{\ordplus}{\mathbin{+}}
\providecommand{\funrestrictionto}[2]{#1 \upharpoonright #2}
\providecommand{\fundefined}{\mathord{\uparrow}}
\providecommand{\Trm}{\mathrm{Trm}}
\providecommand{\OProv}{\operatorname{Prov}}
\providecommand{\OPrf}{\operatorname{Prf}}
\providecommand{\aeq}{\mathrel{=}}
\providecommand{\mClass}[1]{[#1]}
\providecommand{\equivclass}[1]{[#1]}
\providecommand{\pheight}{\operatorname{ht}}
\providecommand{\phy}{\operatorname{ht}}
\providecommand{\Cut}{\mathrm{Cut}}
\providecommand{\QuantRank}{\operatorname{qr}}
\providecommand{\elemequiv}{\mathrel{\equiv}}
\providecommand{\PIso}{\mathrm{PIso}}
\providecommand{\Theory}{\mathrm{Th}}
\providecommand{\supstrict}{\subsetneq}
\providecommand{\trcl}{\operatorname{trcl}}
\providecommand{\cref}[1]{\ref{#1}}
\providecommand{\Discharge}{\mathrm{Dis}}
\providecommand{\Weakening}{\mathrm{Wk}}
\providecommand{\ZFminus}{\mathrm{ZF}^{-}}
\providecommand{\RK}{\mathrm{RK}}
% \Dual is printed both in math mode and (with a bracketed annotation) in running text;
% the math-mode test keeps \mathrm legal in both cases.
\makeatletter
\renewcommand{\Dual}{\@ifnextchar[{\Dual@i}{\ifmmode\mathrm{d}\else$\mathrm{d}$\fi}}
\def\Dual@i[#1]{\ifmmode\mathrm{d}[#1]\else$\mathrm{d}[#1]$\fi}
\makeatother
\providecommand{\RightR}{\mathrm{R}}
\providecommand{\LeftR}{\mathrm{L}}
\providecommand{\Intro}{\mathrm{I}}
\providecommand{\Axiom}{\mathrm{Ax}}
\providecommand{\formula}[1]{#1}
\providecommand{\stared}{\ast}
\providecommand{\oset}[2]{\{\, #1 \mid #2 \,\}}
\providecommand{\setof}[2]{\{\, #1 \mid #2 \,\}}
\providecommand{\pow}[1]{\mathcal{P}(#1)}
\providecommand{\ordpair}[2]{\langle #1, #2 \rangle}
\providecommand{\ordtriple}[3]{\langle #1, #2, #3 \rangle}
\providecommand{\union}{\cup}
\providecommand{\inter}{\cap}
\providecommand{\seq}[1]{\langle #1 \rangle}
\providecommand{\SuchThat}{\mid}
\providecommand{\where}{\mid}
\providecommand{\TVar}{\mathrm{TVar}}
\providecommand{\Bool}{\mathrm{Bool}}
\providecommand{\Val}{\operatorname{Val}}
\providecommand{\FV}{\operatorname{FV}}
\providecommand{\TL}[1]{\mathrm{TL}}
\providecommand{\xstared}{\ast}

% ===================== OLP compatibility layer (5) =====================
% --- fifth compatibility block: rescues the quarantined OLP candidates ---
% The proof-tree and mode fixes live in block 1 (where those macros are first defined); this block
% only adds notation that the quarantined spans still needed and that no earlier block covered.
\providecommand{\ident}{\operatorname{id}}
\providecommand{\setrank}{\operatorname{rank}}
\providecommand{\cardfont}{\mathfrak{c}}
\providecommand{\Complement}{\overline}
\providecommand{\QR}{\mathrm{QR}}
\providecommand{\becd}{\ensuremath{\ulcorner}}
\providecommand{\disjointsum}{\mathbin{\oplus}}
\providecommand{\cardeq}[2]{#1 = #2}
\providecommand{\xred}{\mathrel{\rightarrow}}
\providecommand{\cardless}{<}
\providecommand{\cardplus}{\mathbin{+}}
\providecommand{\cardtimes}{\mathbin{\cdot}}
\providecommand{\cardexpo}{\mathbin{\hat{\ }}}
\providecommand{\cardneq}{\neq}
\providecommand{\cardle}{\leq}
\providecommand{\cardnless}{\not<}
\begin{document}

% ==== OLP file: content/set-theory/card-arithmetic/card-arithmetic.tex ====
% Part: set-theory
% Chapter: card-arithmetic



\olchapter{sth}{card-arithmetic}{Cardinal Arithmetic}

\olimport{opps}
\olimport{simp}
\olimport{expotough}
\olimport{ch}
\olimport{fix}

\OLEndChapterHook

% ==== OLP file: content/set-theory/card-arithmetic/ch.tex ====


\olfileid{sth}{card-arithmetic}{ch}

\olsection{The Continuum Hypothesis}

The previous result hints (correctly) that cardinal exponentiation
would be quite \emph{easy}, if infinite cardinals are guaranteed to
``play straightforwardly'' with powers of $2$, i.e., (by
\olref[opps]{lem:SizePowerset2Exp}) with taking powersets. But we
cannot assume that infinite cardinals \emph{do} play straightforwardly 
powersets. 

To start unpacking this, we introduce some nice notation.

\begin{defn}
Where $\cardsucc{\cardfont{a}}$ is the least cardinal strictly greater
than $\cardfont{a}$, we define two infinite sequences:
\begin{align*}
	\aleph_{0} &\defis \omega & 		
	\beth_{0} &\defis \omega\\
	\aleph_{\alpha \ordplus 1} &\defis \cardsucc{(\aleph_{\alpha})} &
	\beth_{\alpha+1} &\defis \cardexpo{2}{\beth_{\alpha}}\\
	\aleph_{\alpha} &\defis \bigcup_{\beta< \alpha} \aleph_{\beta} &
	\beth_{\alpha} &\defis \bigcup_{\beta < \alpha}\beth_{\beta} & \text{when $\alpha$ is a limit ordinal}.
	\end{align*}
\end{defn}

The definition of $\cardsucc{\cardfont{a}}$ is in order, since
\olref[cardinals][classing]{lem:NoLargestCardinal} tells us that, for each
cardinal $\cardfont{a}$, there is some cardinal greater than
$\cardfont{a}$, and Transfinite Induction guarantees that there is a
\emph{least} cardinal greater than $\cardfont{a}$. The rest of the
definition of $\cardfont{a}$ is provided by transfinite recursion. 

Cantor introduced this ``$\aleph$'' notation; this is \emph{aleph},
the first letter in the Hebrew alphabet and the first letter in the
Hebrew word for ``infinite''. Peirce introduced the ``$\beth$''
notation; this is \emph{beth}, which is the second letter in the
Hebrew alphabet.\footnote{Peirce used this notation in a letter to
Cantor of December 1900. Unfortunately, Peirce also gave a bad
argument there that $\beth_\alpha$ does not exist for $\alpha \geq
\omega$.} Now, these notations provide us with infinite cardinals.

\begin{prop}
$\aleph_\alpha$ and $\beth_\alpha$ are cardinals, for every
ordinal $\alpha$. 
\end{prop}

\begin{proof}
Both results hold by a simple transfinite induction. $\aleph_0 =
\beth_0 = \omega$ is a cardinal by
\olref[cardinals][classing]{omegaisacardinal}. Assuming $\aleph_\alpha$ and
$\beth_\alpha$ are both cardinals, $\aleph_{\alpha+1}$ and
$\beth_{\alpha+1}$ are explicitly defined as cardinals. And the union
of a set of cardinals is a cardinal, by
\olref[cardinals][classing]{unioncardinalscardinal}.
\end{proof}
\noindent
Moreover, every infinite cardinal is an $\aleph$:

\begin{prop}
If $\cardfont{a}$ is an infinite cardinal, then $\cardfont{a} =
\aleph_\gamma$ for some unique $\gamma$.
\end{prop}

\begin{proof}
By transfinite induction on cardinals. For induction, suppose that if
$\cardfont{b} < \cardfont{a}$ then $\cardfont{b} =
\aleph_{\gamma_\cardfont{b}}$. If $\cardfont{a} =
\cardsucc{\cardfont{b}}$ for some $\cardfont{b}$, then $\cardfont{a} =
\cardsucc{(\aleph_{\gamma_\cardfont{b}})}=
\aleph_{\gamma_\cardfont{b}+1}$. If $\cardfont{a}$ is not the
successor of any cardinal, then since cardinals are ordinals
$\cardfont{a} = \bigcup_{\cardfont{b} < \cardfont{a}} \cardfont{b} =
\bigcup_{\cardfont{b} < \cardfont{a}}{\aleph_{\gamma_\cardfont{b}}}$,
so $\cardfont{a} = \aleph_\gamma$ where $\gamma =
\bigcup_{\cardfont{b} < \cardfont{a}}\gamma_\cardfont{b}$. 
\end{proof}

Since every infinite cardinal is an $\aleph$, this prompts us to ask:
is every infinite cardinal a~$\beth$? Certainly if that \emph{were}
the case, then the infinite cardinals would ``play straightforwardly''
with the operation of taking powersets. Indeed, we would have the
following:

\begin{defish}
\emph{Generalized Continuum Hypothesis} (GCH). $\aleph_\alpha  = \beth_\alpha$, for all $\alpha$. 
\end{defish}

Moreover, if GCH held, then we could make some considerable
simplifications with cardinal exponentiation. In particular, we could
show that when $\cardfont{b} < \cardfont{a}$, the value of
$\cardexpo{\cardfont{a}}{\cardfont{b}}$ is trapped by
$\cardfont{a}\leq \cardexpo{\cardfont{a}}{\cardfont{b}} \leq
\cardsucc{\cardfont{a}}$. We could then go on to give precise
conditions which determine which of the two possibilities obtains
(i.e., whether $\cardfont{a} = \cardexpo{\cardfont{a}}{\cardfont{b}}$
or $\cardexpo{\cardfont{a}}{\cardfont{b}} =
\cardsucc{\cardfont{a}}$).\footnote{The condition is dictated by
\emph{cofinality}.}

But GCH is a \emph{hypothesis}, not a \emph{theorem}. In fact,
\citet{Godel1938} proved that if $\ZFC$ is consistent, then so is
$\ZFC + \text{GCH}$. But it later turned out that we can equally add
$\lnot$GCH to $\ZFC$. Indeed, consider the simplest non-trivial
\emph{instance} of GCH, namely: 

\begin{defish}
\emph{Continuum Hypothesis} (CH). $\aleph_1 = \beth_1$. 
\end{defish}

\citet{Cohen1963} proved that if $\ZFC$ is consistent then so is $\ZFC
+ \lnot\text{CH}$. So the Continuum Hypothesis is independent from $\ZFC$.

The Continuum Hypothesis is so-called, since ``the continuum'' is
another name for the real line, $\Real$.
\olref[opps]{continuumis2aleph0} tells us that $\card{\Real} =
\beth_1$. So the Continuum Hypothesis states that there is no cardinal
between the cardinality of the natural numbers, $\aleph_0 = \beth_0$,
and the cardinality of the continuum, $\beth_1$.

Given the \emph{independence} of (G)CH from $\ZFC$, what should say
about their \emph{truth}? Well, there is \emph{much} to say. Indeed,
and much fertile recent work in set theory has been directed at
investigating these issues. But two very quick points are certainly worth
emphasising. 

First: it does not \emph{immediately} follow from these formal
independence results that either GCH or CH is \emph{indeterminate} in
truth value. After all, maybe we just need to add more axioms, which
strike us as natural, and which will settle the question one way or
another. G\"odel himself suggested that this was the right response. 

Second: the independence of CH from $\ZFC$ is certainly
\emph{striking}, but it is certainly not \emph{incredible} (in the
literal sense). The point is simply that, for all $\ZFC$ tells us,
moving from cardinals to their successors may involve a less blunt
tool than simply taking powersets.
 %The operation of taking powersets moves us from one stage of the
 %hierarchy to its successor stage (from $V_{\alpha}$ to
 %$V_{\alpha+1}$). 

With those two observations made, if you want to know more, you will
now have to turn to the various philosophers and mathematicians with
horses in the race.\footnote{Though you might want to start by reading \citet[\S15.6]{Potter2004}.}

% ==== OLP file: content/set-theory/card-arithmetic/expotough.tex ====

\olfileid{sth}{card-arithmetic}{expotough}

\olsection[Some Simplifications]{Some Simplification with Cardinal Exponentiation}

Whilst defining $\canonord$ was a little involved, the upshot is a
useful result concerning cardinal addition and multiplication,
\olref[simp]{cardplustimesmax}. Transfinite exponentiation, however,
cannot be simplified so straightforwardly. To explain why, we start
with a result which extends a familiar pattern from the finitary case
(though its proof is at a high level of abstraction):

\begin{prop}\ollabel{simplecardexpo}
$\cardexpo{\cardfont{a}}{\cardfont{b} \cardplus \cardfont{c}} =
\cardexpo{\cardfont{a}}{\cardfont{b}} \cardtimes
\cardexpo{\cardfont{a}}{\cardfont{c}}$ and
$\cardexpo{(\cardexpo{\cardfont{a}}{\cardfont{b}})}{\cardfont{c}} =
\cardexpo{\cardfont{a}}{\cardfont{b} \cardtimes \cardfont{c}}$, for
any cardinals $\cardfont{a}, \cardfont{b}, \cardfont{c}$.
\end{prop}

\begin{proof}
For the first claim, consider a function $f \colon
(\cardfont{b}\disjointsum\cardfont{c}) \to \cardfont{a}$. Now ``split
this'', by defining $f_\cardfont{b}(\beta) = f(\beta, 0)$ for each
$\beta \in \cardfont{b}$, and $f_\cardfont{c}(\gamma) = f(\gamma, 1)$
for each $\gamma \in \cardfont{c}$. The map $f \mapsto
(f_{\cardfont{b}} \times f_\cardfont{c})$ is !!a{bijection}
$\funfromto{\cardfont{b} \disjointsum \cardfont{c}}{\cardfont{a}} \to
(\funfromto{\cardfont{b}}{\cardfont{a}} \times
\funfromto{\cardfont{c}}{\cardfont{a}})$. 

For the second claim, consider a function $f \colon \cardfont{c} \to
(\funfromto{\cardfont{b}}{\cardfont{a}})$; so for each $\gamma \in
\cardfont{c}$ we have some function $f(\gamma) \colon \cardfont{b} \to
\cardfont{a}$. Now define $f^*(\beta, \gamma) = (f(\gamma))(\beta)$
for each $\tuple{\beta, \gamma} \in \cardfont{b} \times \cardfont{c}$.
The map $f \mapsto f^*$ is !!a{bijection}
$\funfromto{\cardfont{c}}{(\funfromto{\cardfont{b}}{\cardfont{a}})}
\to \funfromto{\cardfont{b} \cardtimes \cardfont{c}}{\cardfont{a}}$. 
\end{proof}

Now, what we would \emph{like} is an easy way to compute
$\cardexpo{\cardfont{a}}{\cardfont{b}}$ when we are dealing with
infinite cardinals. Here is a nice step in this direction:

\begin{prop}\ollabel{cardexpo2reduct}
If $2 \leq \cardfont{a} \leq \cardfont{b}$ and $\cardfont{b}$ is
infinite, then $\cardexpo{\cardfont{a}}{\cardfont{b}} =
\cardexpo{2}{\cardfont{b}}$
\end{prop}

\begin{proof}
\begin{align*}
	\cardexpo{2}{\cardfont{b}} &\leq 
	\cardexpo{\cardfont{a}}{\cardfont{b}}\text{, as $2 \leq \cardfont{a}$}\\
	&\leq \cardexpo{(2^\cardfont{a})}{\cardfont{b}}
	\text{, by \olref[opps]{lem:SizePowerset2Exp}}\\
	&= \cardexpo{2}{\cardfont{a} \cardtimes \cardfont{b}}
	\text{, by \olref{simplecardexpo}} \\
	&= \cardexpo{2}{\cardfont{b}}
	\text{, by \olref[simp]{cardplustimesmax}}
\end{align*}
\end{proof}

We should not really expect to be able to simplify this any further,
since $\cardfont{b} < \cardexpo{2}{\cardfont{b}}$ by
\olref[card-arithmetic][opps]{lem:SizePowerset2Exp}.
However, this does not tell us what to say about
$\cardexpo{\cardfont{a}}{\cardfont{b}}$ when $\cardfont{b} <
\cardfont{a}$. Of course, if $\cardfont{b}$ is \emph{finite}, we know
what to do.

\begin{prop}
If $\cardfont{a}$ is infinite and $n \in \omega$ then
$\cardexpo{\cardfont{a}}{n} = \cardfont{a}$
\end{prop}

\begin{proof}
$\cardexpo{\cardfont{a}}{n} = \cardfont{a} \cardtimes \cardfont{a}
\cardtimes \ldots \cardtimes \cardfont{a} = \cardfont{a}$, by \olref[simp]{cardplustimesmax}.
\end{proof}
\noindent 
Additionally, in some other cases, we can control the size of
$\cardexpo{\cardfont{a}}{\cardfont{b}}$:

\begin{prop}
If $2 \leq \cardfont{b} < \cardfont{a} \leq
\cardexpo{2}{\cardfont{b}}$ and $\cardfont{b}$ is infinite, then
$\cardexpo{\cardfont{a}}{\cardfont{b}} = \cardexpo{2}{\cardfont{b}}$
\end{prop}

\begin{proof}
$\cardexpo{2}{\cardfont{b}}\leq \cardexpo{\cardfont{a}}{\cardfont{b}}
\leq \cardexpo{(\cardexpo{2}{\cardfont{b}})}{\cardfont{b}} =
\cardexpo{2}{\cardfont{b}\cardtimes\cardfont{b}} =
\cardexpo{2}{\cardfont{b}}$, reasoning as in \olref{cardexpo2reduct}.
\end{proof}
\noindent 
But, beyond this point, things become rather more subtle.

% ==== OLP file: content/set-theory/card-arithmetic/fix.tex ====


\olfileid{sth}{card-arithmetic}{fix}
\olsection{$\aleph$-Fixed Points}

In \olref[spine][]{chap}, we suggested that Replacement stands in need
of justification, because it forces the hierarchy to be rather tall.
Having done some cardinal arithmetic, we can give a little
illustration of the height of the hierarchy. 

Evidently $0 < \aleph_0$, and $1 < \aleph_1$, and $2 < \aleph_2$\ldots
and, indeed, the difference in size only gets \emph{bigger} with every
step. So it is tempting to conjecture that $\kappa< \aleph_\kappa$
for every ordinal $\kappa$. 

But this conjecture is \emph{false}, given $\ZFC$. In fact, we can
prove that there are \emph{$\aleph$-fixed-points}, i.e.,
cardinals $\kappa$ such that $\kappa=\aleph_\kappa$. 

\begin{prop}\ollabel{alephfixed}
There is an $\aleph$-fixed-point.
\end{prop}

\begin{proof}
Using recursion, define:
\begin{align*}
	\kappa_0 &= 0\\
	\kappa_{n+1} &= \aleph_{\kappa_n}\\
	\kappa&= \bigcup_{n < \omega}\kappa_n
\end{align*}
Now $\kappa$ is a cardinal by
\olref[cardinals][classing]{unioncardinalscardinal}. But now:
\[
	\kappa= \bigcup_{n < \omega} \kappa_{n+1} = 
	\bigcup_{n < \omega}\aleph_{\kappa_n} = 
	\bigcup_{\alpha < \kappa}\aleph_\alpha = \aleph_\kappa
\]
%	By construction, $\kappa$ is the least cardinal greater than each
%	$\kappa_n$. So $\aleph_\kappa$ is the least cardinal greater than
%	each $\aleph_{\kappa_n}$, and hence greater than each
%	$\kappa_{n+1}$. But equally $\kappa$ is the least cardinal greater
%	than each $\kappa_{n+1} = \aleph_{\kappa_n}$. So $\kappa=
%	\aleph_\kappa$.
\end{proof}

Boolos once wrote an article about exactly the $\aleph$-fixed-point we
just constructed. After noting the existence of $\kappa$, at the start
of his article, he said:
\begin{quote}
	[$\kappa$ is] a \emph{pretty big} number, by the lights of those
	with no previous exposure to  set theory,  so big, it seems to me,
	that  it calls into question the truth of any theory, one of whose
	assertions is the claim that there are at least $\kappa$ objects.
	\citep[p.~257]{Boolos2000}
\end{quote}
And he ultimately concluded his paper by asking:
\begin{quote}
	[do] we  suspect that,  however  it  may  have  been  at  the
	beginning  of  the  story,  by  the  time  we have come thus  far
	the wheels  are  spinning  and  we  are no longer  listening  to
	a description  of  anything  that  is the case?
	\citep[p.~268]{Boolos2000}
\end{quote}
If we have, indeed, outrun ``anything that is the case'', then we must
point the finger of blame directly at Replacement. For it is this
axiom which allows our proof to work. In which case, one assumes,
Boolos would need to revisit the claim he made, a few decades earlier,
that Replacement has ``no undesirable'' consequences (see
\olref[replacement][extrinsic]{sec}).

But is the existence of $\kappa$ so bad? It might help, here, to
consider Russell's \emph{Tristram Shandy paradox}. Tristram Shandy
documents his life in his diary, but it takes him a year to record a
single day. With every passing year, Tristram falls further and
further behind: after one year, he has recorded only one day, and has
lived 364 days unrecorded days; after two years, he has only recorded
two days, and has lived 728 unrecorded days; after three years, he has
only recorded three days, and lived 1092 unrecorded
days \dots\footnote{Forgetting about leap years.} Still, if Tristram
is \emph{immortal}, Tristram will manage to record every day, for he
will record the $n$th day on the $n$th year of his life. And so, ``at
the end of time'', Tristram will have a complete diary. 

Now: why is this so different from the thought that $\alpha$ is
smaller than $\aleph_\alpha$---and indeed, increasingly, desperately
smaller---up until $\kappa$, at which point, we catch up, and $\kappa
= \aleph_\kappa$?

Setting that aside, and assuming we accept $\ZFC$, let's close with a
little more fun concerning fixed-point constructions. The next three
results establish, intuitively, that there is a (non-trivial) point at
which the hierarchy is as wide as it is tall:

\begin{prop}\ollabel{bethfixed}
There is a $\beth$-fixed-point, i.e., a $\kappa$ such that $\kappa=
\beth_\kappa$.
\end{prop}

\begin{proof}
As in \olref{alephfixed}, using ``$\beth$'' in place of ``$\aleph$''. 
\end{proof}

\begin{prop}\ollabel{stagesize}
$\card{V_{\omega+\alpha}} = \beth_{\alpha}$. If $\omega \ordtimes
\omega \leq \alpha$, then $\card{V_\alpha} = \beth_\alpha$.
\end{prop}

\begin{proof}
The first claim holds by a simple transfinite induction. The second
claim follows, since if $\omega \ordtimes \omega \leq \alpha$ then
$\omega + \alpha = \alpha$. To establish this, we use facts about
ordinal arithmetic from \olref[ord-arithmetic][]{chap}. First note
that $\omega \ordtimes \omega = \omega \ordtimes (1 \ordplus \omega) =
(\omega  \ordtimes 1) \ordplus (\omega\ordtimes\omega) = \omega
\ordplus (\omega \ordtimes \omega)$. Now if $\omega \ordtimes \omega
\leq \alpha$, i.e., $\alpha = (\omega\ordtimes\omega) \ordplus \beta$
for some $\beta$, then $\omega \ordplus \alpha = \omega \ordplus
((\omega \ordtimes \omega) \ordplus \beta) = (\omega \ordplus (\omega
\ordtimes \omega)) \ordplus \beta = (\omega \ordtimes \omega) \ordplus
\beta = \alpha$. 
\end{proof}

\begin{cor}
There is a $\kappa$ such that $\card{V_\kappa} = \kappa$.
\end{cor}

\begin{proof}
Let $\kappa$ be a $\beth$-fixed point, as given by \olref{bethfixed}.
Clearly $\omega \ordtimes \omega < \kappa$. So $\card{V_\kappa} =
\beth_\kappa= \kappa$ by \olref{stagesize}.
\end{proof}

There are as many stages beneath $V_\kappa$ as there are !!{element}s
of $V_\kappa$. Intuitively, then, $V_\kappa$ is as wide as it is tall.
This is very Tristram-Shandy-esque: we move from one stage to the next
by taking \emph{powersets}, thereby making our hierarchy \emph{much}
bigger with each step. But, ``in the end'', i.e., at stage $\kappa$,
the hierarchy's width catches up with its height. 

One might ask: \emph{How often does the hierarchy's width match its
height?} The answer is: \emph{As often as there are ordinals.} But
this needs a little explanation. 

We define a term $\tau$ as follows. For any $A$, let:
\begin{align*}
	\tau_0(A) & \defis \card{A}\\
	\tau_{n+1}(A) & \defis \beth_{\tau_n(A)}\\
	\tau(A) & \defis \bigcup_{n < \omega}\tau_n(A)
\intertext{As in \olref{bethfixed}, $\tau(A)$ is a
$\beth$-fixed point for any $A$, and trivially $\card{A} < \tau(A)$.
So now consider this recursive definition:}
	W_0 &\defis 0\\
	W_{\alpha + 1} & \defis \tau(W_\alpha)\\
	W_\alpha & \defis \bigcup_{\beta < \alpha} W_\beta
	\text{, when $\alpha$ is a limit}
\end{align*}
The construction is defined for all ordinals. Intuitively, then,
$W$ is ``!!a{injection}'' from the ordinals to $\beth$-fixed points.
And, exactly as before, 
$V_{W_\alpha}$ is as wide as it is tall, for any $\alpha$.

% ==== OLP file: content/set-theory/card-arithmetic/opps.tex ====


\olfileid{sth}{card-arithmetic}{opps}
\olsection{Defining the Basic Operations}

Since we do not need to keep track of order, cardinal arithmetic is
rather easier to define than ordinal arithmetic. We will define
addition, multiplication, and exponentiation simultaneously. 

\begin{defn}
When $\cardfont{a}$ and $\cardfont{b}$ are cardinals:
\begin{align*}
	\cardfont{a} \cardplus \cardfont{b} &\defis 
	\card{\cardfont{a} \disjointsum \cardfont{b}}\\
	\cardfont{a} \cardtimes \cardfont{b} &\defis 
	\card{\cardfont{a} \times \cardfont{b}}\\
	\cardexpo{\cardfont{a}}{\cardfont{b}} &\defis 
	\card{\funfromto{\cardfont{b}}{\cardfont{a}}}
\end{align*}
where $\funfromto{X}{Y} = \Setabs{f}{f \text{ is a function } X \to Y}$.
(It is easy to show that $\funfromto{X}{Y}$ exists for any sets $X$
and $Y$; we leave this as an exercise.) 
\end{defn}

\begin{prob}
Prove in $\Zminus$ that $\funfromto{X}{Y}$ exists for any sets $X$
and~$Y$. Working in $\ZF$, compute $\setrank{\funfromto{X}{Y}}$ from
$\setrank{X}$ and $\setrank{Y}$, in the manner of
\olref[sth][ord-arithmetic][using-addition]{rankcomputation}. 
\end{prob}

It might help to explain this definition. Concerning addition: this
uses the notion of disjoint sum, $\disjointsum$, as defined in
\olref[ord-arithmetic][add]{defdissum}; and it is easy
to see that this definition gives the right verdict for finite cases.
Concerning multiplication: \olref[sfr][set][pai]{cardnmprod} tells us
that if $A$ has $n$ members and $B$ has $m$ members then $A \times B$
has $n \cdot m$ members, so our definition simply generalises the idea
to transfinite multiplication. Exponentiation is similar: we are
simply generalising the thought from the finite to the transfinite.
Indeed, in certain ways, transfinite cardinal arithmetic looks much
more like ``ordinary'' arithmetic than does transfinite ordinal
arithmetic:

\begin{prop}\ollabel{cardplustimescommute}
$\cardplus$ and $\cardtimes$ are commutative and associative. 
\end{prop}

\begin{proof}
For commutativity, by
\olref[cardinals][cardsasords]{lem:CardinalsBehaveRight} it
suffices to observe that $\cardeq{(\cardfont{a} \disjointsum
\cardfont{b})}{(\cardfont{b} \disjointsum \cardfont{a})}$ and
$\cardeq{(\cardfont{a} \times \cardfont{b})}{(\cardfont{b} \times
\cardfont{a})}$. We leave associativity as an exercise.
\end{proof}

\begin{prob}
Prove that $\cardplus$ and $\cardtimes$ are associative.
\end{prob}

\begin{prop}
$A$ is infinite iff $\card{A} \cardplus 1 = 1 \cardplus \card{A} = \card{A}$.
\end{prop}

\begin{proof}
As in
\olref[cardinals][classing]{generalinfinitycharacter}, from
\olref[ord-arithmetic][using-addition]{ordinfinitycharacter} and
\olref[cardinals][cardsasords]{lem:CardinalsBehaveRight}. 
\end{proof}

This explains why we need to use different symbols for ordinal versus
cardinal addition/multiplication: these are genuinely \emph{different}
operations. This next pair of results shows that ordinal versus
cardinal exponentiation are also different operations. (Recall that
\olref[z][infinity-again]{defnomega} entails that $2 = \{0,
1\}$):

\begin{lem}\ollabel{lem:SizePowerset2Exp}
$\card{\Pow{A}} = \cardexpo{2}{\card{A}}$, for any $A$.
\end{lem}

\begin{proof}
For each subset $B \subseteq A$, let $\chi_B \in \funfromto{A}{2}$ be given by:
\begin{align*}
	\chi_{B}(x) &\defis
	\begin{cases}
		1 & \text{if }x\in B\\
		0 & \text{otherwise.}
	\end{cases}
\end{align*}
Now let $f(B) = \chi_B$; this defines !!a{bijection} $f \colon \Pow{A}
\to \funfromto{A}{2}$. So $\cardeq{\Pow{A}}{\funfromto{A}{2}}$. Hence
$\cardeq{\Pow{A}}{\funfromto{\card{A}}{2}}$, so that
$\card{\Pow{A}} = \card{\funfromto{\card{A}}{2}} =
2^{\card{A}}$.
\end{proof}

This snappy proof essentially subsumes the discussion of
\olref[sfr][siz][red-alt]{sec}. There, we showed how to ``reduce'' the
uncountability of $\Pow{\omega}$ to the uncountability of the set of
infinite binary strings, $\Bin^\omega$. In effect, $\Bin^{\omega}$ is
just $\funfromto{\omega}{2}$; and the preceding proof showed that the
reasoning we went through in \olref[sfr][siz][red-alt]{sec} will go
through using any set~$A$ in place of~$\omega$. The result also yields
a quick fact about cardinal exponentiation:

\begin{cor}\ollabel{cantorcor}
$\cardfont{a} < \cardexpo{2}{\cardfont{a}}$ for any cardinal~$\cardfont{a}$.
\end{cor}

\begin{proof}
From Cantor's Theorem (\olref[sfr][siz][car]{thm:cantor}) and
\olref{lem:SizePowerset2Exp}.
\end{proof}
\noindent
So $\omega < \cardexpo{2}{\omega}$. But note: this is a result about
\emph{cardinal} exponentiation. It should be contrasted with
\emph{ordinal} exponentiation, since in the latter case $\omega =
\ordexpo{2}{\omega}$ (see \olref[ord-arithmetic][expo]{sec}).

Whilst we are on the topic of cardinal exponentiation, we can also be
a bit more precise about the ``way'' in which $\Real$ is
!!{nonenumerable}.

\begin{thm}\ollabel{continuumis2aleph0}
$\card{\Real} = \cardexpo{2}{\omega}$
\end{thm}

\begin{proof}[Proof skeleton]
There are plenty of ways to prove this. The most straightforward is to
argue that $\cardle{\Pow{\omega}}{\Real}$ and
$\cardle{\Real}{\Pow{\omega}}$, and then use Schr\"oder-Bernstein to
infer that $\cardeq{\Real}{\Pow{\omega}}$, and
\olref[card-arithmetic][opps]{lem:SizePowerset2Exp} to infer
that $\card{\Real} = \cardexpo{2}{\omega}$. We leave it as an
(illuminating) exercise to define injections $f \colon
\Pow{\omega} \to \Real$ and $g \colon \Real \to \Pow{\omega}$.
\end{proof}

\begin{prob}
Complete the proof of
\olref[sth][card-arithmetic][opps]{continuumis2aleph0}, by
showing that $\cardle{\Pow{\omega}}{\Real}$ and
$\cardle{\Real}{\Pow{\omega}}$.
\end{prob}

% ==== OLP file: content/set-theory/card-arithmetic/simp.tex ====

\olfileid{sth}{card-arithmetic}{simp}

\olsection{Simplifying Addition and Multiplication}

It turns out that transfinite cardinal addition and multiplication is
\emph{extremely} easy. This follows from the fact that cardinals are
(certain) ordinals, and so well-ordered, and so can be manipulated in
a certain way. Showing this, though, is \emph{not} so easy. To start,
we need a tricksy definition:

\begin{defn}
We define a \emph{canonical ordering}, $\canonord$, on pairs of
ordinals, by stipulating that $\tuple{\alpha_1, \alpha_2} \canonord
\tuple{\beta_1, \beta_2}$ iff either:
\begin{enumerate}
	\item $\max(\alpha_1, \alpha_2) < \max(\beta_1, \beta_2)$; or
	\item $\max(\alpha_1, \alpha_2) = \max(\beta_1, \beta_2)$ and
	$\alpha_1 < \beta_1$; or
	\item $\max(\alpha_1, \alpha_2) = \max(\beta_1, \beta_2)$ and
	$\alpha_1 = \beta_1$ and $\alpha_2 < \beta_2$
\end{enumerate}
\end{defn}

\begin{lem}
$\tuple{\alpha \times \alpha, \canonord}$ is a well-order, for any
ordinal $\alpha$.
\end{lem}

\begin{proof}
Evidently $\canonord$ is connected on $\alpha \times \alpha$. For
suppose that neither $\tuple{\alpha_1, \alpha_2}$ nor $\tuple{\beta_1,
\beta_2}$ is $\canonord$-less than the other. Then $\max(\alpha_1,
\alpha_2) = \max(\beta_1, \beta_2)$ and $\alpha_1 = \beta_1$ and
$\alpha_2 = \beta_2$, so that $\tuple{\alpha_1, \alpha_2} =
\tuple{\beta_1,  \beta_2}$.

To show well-ordering, let $X \subseteq \alpha\times\alpha$ be
non-empty. Since $\alpha$ is an ordinal, some $\delta$ is the least
member of $\Setabs{\max(\gamma_1, \gamma_2)}{\tuple{\gamma_1,
\gamma_2} \in X}$. Now discard all pairs from
$\Setabs{\tuple{\gamma_1,\gamma_2} \in X}{\max(\gamma_1, \gamma_2) =
\delta}$ except those with least first coordinate; from among these,
the pair with least second coordinate is the $\canonord$-least element
of $X$.
\end{proof}
\noindent
Now for a teensy, simple observation:

\begin{prop}\ollabel{simplecardproduct}
If $\cardeq{\alpha}{\beta}$, then $\cardeq{\alpha \times \alpha}{\beta
\times \beta}$. 
\end{prop}

\begin{proof}
Just let $f \colon \alpha \to \beta$ induce $\tuple{\gamma_1,
\gamma_2} \mapsto \tuple{f(\gamma_1), f(\gamma_2)}$.
\end{proof}
\noindent
And now we will put all this to work, in proving a crucial lemma:
\begin{lem}\ollabel{alphatimesalpha}
$\cardeq{\alpha}{\alpha \times \alpha}$, for any infinite ordinal
$\alpha$
\end{lem}

\begin{proof}
For reductio, let $\alpha$ be the least infinite ordinal for which
this is false. \olref[sfr][siz][zigzag]{natsquaredenumerable} shows
that $\cardeq{\omega}{\omega\times\omega}$, so $\omega \in \alpha$.
Moreover, $\alpha$ is a cardinal: suppose otherwise, for reductio;
then $\card{\alpha} \in \alpha$, so that
$\cardeq{\card{\alpha}}{\card{\alpha} \times \card{\alpha}}$, by
hypothesis; and $\cardeq{\card{\alpha}}{\alpha}$ by definition; so
that $\cardeq{\alpha}{\alpha\times\alpha}$ by
\olref{simplecardproduct}. 

Now, for each $\tuple{\gamma_1, \gamma_2} \in \alpha \times \alpha$,
consider the segment:
\begin{align*}
	\text{Seg}(\gamma_1, \gamma_2) &= \Setabs{\tuple{\delta_1, \delta_2} \in \alpha \times \alpha}{\tuple{\delta_1, \delta_2} \canonord \tuple{\gamma_1, \gamma_2}}
\end{align*}
Letting $\gamma = \max(\gamma_1, \gamma_2)$, note that $\tuple{\gamma_1, \gamma_2} \canonord \tuple{\gamma+1, \gamma + 1}$. So, when $\gamma$ is infinite, observe:
\begin{align*}
	\text{Seg}(\gamma_1, \gamma_2) & 
	\precsim ((\gamma \ordplus 1)\ordtimes (\gamma \ordplus 1))\\
	&\approx (\gamma \ordtimes \gamma)
	\text{, by \olref[ord-arithmetic][using-addition]{ordinfinitycharacter} and 
	\olref{simplecardproduct}}\\
	&\approx \gamma \text{, by the induction hypothesis}\\
	& \prec \alpha\text{, since $\alpha$ is a cardinal}
\end{align*}
So $\ordtype{\alpha\times \alpha, \canonord} \leq \alpha$, and hence
$\cardle{\alpha \times \alpha}{\alpha}$. Since of course
$\cardle{\alpha}{\alpha \times \alpha}$, the result follows by
Schr\"oder-Bernstein. 
\end{proof}

Finally, we get to our simplifying result:

\begin{thm}\ollabel{cardplustimesmax}
If $\cardfont{a}, \cardfont{b}$ are infinite cardinals, then:
\[
	\cardfont{a}
\cardtimes \cardfont{b} = \cardfont{a} \cardplus \cardfont{b} =
\text{max}(\cardfont{a}, \cardfont{b}).
\]
\end{thm}

\begin{proof}
Without loss of generality, suppose $\cardfont{a} = \max(\cardfont{a},
\cardfont{b})$. Then invoking \olref{alphatimesalpha},
$\cardfont{a}\cardtimes\cardfont{a} = \cardfont{a} \leq \cardfont{a}
\cardplus \cardfont{b} \leq \cardfont{a} \cardplus \cardfont{a} \leq
\cardfont{a} \cardtimes \cardfont{a}$. \end{proof}\noindent Similarly,
if $\cardfont{a}$ is infinite, an $\cardfont{a}$-sized union of
$\leq\cardfont{a}$-sized sets has size $\leq\cardfont{a}$:

\begin{prop}\ollabel{kappaunionkappasize}
Let $\cardfont{a}$ be an infinite cardinal. For each ordinal $\beta
\in \cardfont{a}$, let $X_\beta$ be a set with $\card{X_\beta} \leq
\cardfont{a}$. Then $\card{\bigcup_{\beta \in \cardfont{a}} X_\beta}
\leq \cardfont{a}$.
\end{prop}

\begin{proof}
For each $\beta \in \cardfont{a}$, fix !!a{injection} $f_\beta \colon
X_\beta \to \cardfont{a}$.\footnote{How are these ``fixed''? See \olref[sth][choice][countablechoice]{sec}.} Define !!a{injection} $g \colon
\bigcup_{\beta \in \cardfont{a}} X_\beta \to \cardfont{a} \times
\cardfont{a}$ by $g(v) = \tuple{\beta, f_\beta(v)}$, where $v \in
X_\beta$ and $v \notin X_\gamma$ for any $\gamma \in \beta$. Now
$\bigcup_{\beta \in \cardfont{a}} X_\beta \preceq \cardfont{a} \times
\cardfont{a} \approx \cardfont{a}$ by \olref{cardplustimesmax}.
\end{proof}

